% "Where Should the Bits Go? A First-Order Analysis of KV Cache Quantization
%  with Optimal Key--Value Bit Allocation"
% Plan Week 1: stub all section headers + a one-line claim each. Build: pdflatex main.
\documentclass[11pt]{article}

\usepackage{amsmath,amssymb,amsthm}
\usepackage{graphicx}
\usepackage{hyperref}

\newtheorem{theorem}{Theorem}
\newtheorem{lemma}{Lemma}
\newtheorem{corollary}{Corollary}

\title{Where Should the Bits Go? \\
  \large A First-Order Analysis of KV Cache Quantization with
  Optimal Key--Value Bit Allocation}
\author{Nazmul Alam Diptu}
\date{2026}

\begin{document}
\maketitle

\begin{abstract}
% Week 7: write once the result's shape is known.
KV cache quantization schemes consistently give keys more bits than values, but the
allocation is hand-tuned. We derive a first-order perturbation bound for attention
outputs under stochastic uniform quantization that yields a closed-form optimal
key--value bit allocation, predictable per layer and per head from statistics
measurable on a running model.
\end{abstract}

\section{Introduction}          % Claim: the K/V asymmetry is derivable, not tuned.
% Section: Introduction. Draft in Week 7 (abstract + intro once result shape known).
% One-line claim: keys and values are not equally hard to quantize, and the optimal
% bit gap is derivable from measurable statistics rather than tuned.


\section{Theory}                % Claim: value error is linear; key error amplifies through softmax.
% Section: Theory. THE CORE. Value-side bound = Week 1; key-side + corollary = Weeks 2-3.
% Lock this section at the Week 3 gate once the numerical check in tests/test_bound.py passes.

% Lemma (value-side, easy half): attention output o = a·V is a convex combination
% (weights sum to 1), so value perturbations pass through at most linearly.
\begin{lemma}[Value error does not amplify]
% TODO(Week 1): statement + proof.
\end{lemma}

% Lemma (key-side, hard half): a key perturbation is scaled by the query norm, passed
% through the softmax map (Lipschitz), then scaled by the spread of the values.
\begin{lemma}[Key error amplifies through softmax]
% TODO(Week 2-3): softmax-Lipschitz statement + proof. Fallback: looser triangle bound.
\end{lemma}

\begin{theorem}[First-order output-error bound]
% \|õ − o\| ≤ C · (‖q‖/√d) · diam(V) · r_K · 2^{−b_K} + r_V · 2^{−b_V}
% TODO(Week 2-3).
\end{theorem}

\begin{corollary}[Optimal bit allocation]
% b_K − b_V ≈ log2( ‖q‖ · diam(V) / √d ) + log2( r_K / r_V )
% TODO(Week 2-3): Lagrangian on the bound at fixed target error.
\end{corollary}


\section{Experiments}           % Claim: the bound tracks measured error; derived allocation wins at matched bits.
% Section: Experiments. Write observations here AS runs land (Weeks 5-7), not after.
% Claims to support:
%   - predicted b_K − b_V > 0 (keys need more bits) and the gap varies by layer;
%   - the bound tracks measured attention-output error across 2/3/4-bit;
%   - derived allocation beats uniform and fixed-asymmetric at matched average bits.


\section{Related Work}          % Claim: KIVI showed the asymmetry; we predict its magnitude.
% Section: Related Work. Write DIRECTLY here while reading (standing rule #2) — Weeks 3-4.
% One line each: KIVI (asymmetry empirical, we derive it), KVQuant (outlier/layer stats),
% QJL (theorem-first, closest in spirit), PolarQuant (rotation+residual, orthogonal),
% KV Cache Transform Coding (transform-coding view; our step = reverse water-filling).


\section{Limitations and Future Work}
% Claim: first-order only; rotations compose as the natural sequel.

\end{document}
